\documentclass[11pt]{article}
\usepackage[T1]{fontenc}
\usepackage[utf8]{inputenc}


% 1-inch margins
\usepackage[margin=1in]{geometry} 

% Double spacing
\usepackage{setspace}

% set Times New Roman-style font 
%\usepackage{newtxtext}
%\usepackage{newtxmath}
\usepackage[scaled]{helvet}
\renewcommand{\familydefault}{\sfdefault}


\begin{document} 

\title{Practitioner Notes}
\author{\noauthor}
\date{\nodate}
\maketitle

\section{What is already known about this topic}
\begin{itemize}
    \item Educational technologies are widely used to support learning. While learning, students are often permitted to skip problems or otherwise not attempt it, in order to regulate frustration and sustain motivation. 
    \item Individual differences in the use of meta-cognitive strategies result in differences in the use of non-attempt options. 
\end{itemize}

\section{What this paper adds}
\begin{itemize}
    \item 
\end{itemize}

\section{Implications for practice and/or policy}
\begin{itemize}
    \item Educational technologies rely on accurate student models to estimate and personalize student learning. Individual differences in problem-skipping threaten the accuracy of the common unidimensional student model. Educational technologies who offer non-attempt options therefore risk distorting ability estimates, which, in turn, may compromise personalized recommendations. 
    \item In practice, individual tendencies for non-attempting can be accounted for 
\end{itemize}

\end{document}
